% \clearpage
%%%%%%%%%%%%%%%%%%%%%%%%%%%%%%%%%%%%%%%%%%%%%%%%%%%%%%%%%%%
%%% Abdomen                                             %%%
%%%%%%%%%%%%%%%%%%%%%%%%%%%%%%%%%%%%%%%%%%%%%%%%%%%%%%%%%%%
\begin{savequote}
\includegraphics[width=5cm]{./images/operation.jpg}
\end{savequote}
\chapter
[\CO{[ABD]} Abdominale Erkrankungen]
{\CC{[ABD]} Abdominale Erkrankungen}
\ChapterIntro
{%
\Captionof{figure}{Oft die
    Abteilung der Wahl für einen abdominellen Notfall: Die
    Chirurgie. Aber auch die Abteilungen für
    Kinderheilkunde, Kinderchirurgie, Gynäkologie (!) oder
    die Interne können manchmal passende Ziele sein.
    \SourcePicture{Gabriel}}
\SlideCompensationNoParskip{1}
Der Bauch ist eine echte Herausforderung. Er beinhaltet
viele unterschiedliche Organe und Strukturen welche weh tun
und erkranken können. Viele dieser, oft auch mit starken
Schmerzen verbundenen Störungen sind nicht übermäßig
schlimm, einige sind jedoch sehr gefährlich oder können
rasch gefährlich werden.


}
{%
\begin{description}
  \item[Maintainer:] Sebastian Gabriel
  \item[Autoren:] Diverse
  \item[Reviewer:] Standard-Reviewprozess
  \item[Version:] {\DocVersion} ({\DocVersionInfo})
  \item[SHA1:] {\DocGitCommit}
\end{description}

{\TxTeilkapitelAASS}
}


\clearpage

\section{Allgemeines}

\Layout{0}{Anamnese und Untersuchung}{%
Es ist daher wichtig
typische Symptome zu erheben und Alarmzeichen sofort zu
erkennen.
Wesentlich ist bei allen unklaren Bauchbeschwerden die
Erhebung einer exakten \I{Stuhlanamnese}:

\begin{itemize}
  \item Auffälligkeiten? (Geruch, Farbe, Konsistenz, Frequenz)
  \item Blutbeimengungen
  \item Wann erfolgte der letzte Stuhlgang?
  \item Winde?
  \item Genaue Beschreibung von Harn und eventuell Erbrochenem.
  \item Bei Frauen im gebärfähigem Alter Regelanamnese.
\end{itemize}

Der Patient muss bis zur endgültigen Klärung der Situation
\EE{nüchtern} belassen werden!
Wenn es der Patient toleriert, muss eine Betrachtung und
Abtastung des \so{unbekleideten} Abdomens beim liegenden
Patienten erfolgen wie unter
\ref{topic:untersuchung-abdomen-abtasten} beschrieben.
}{%
\begin{itemize}
  \item Stuhl
  \begin{itemize}
    \item Wann letzter?
    \item Auffälligkeiten
    \item Farbe, Konsistenz
  \end{itemize}
  \item Erbrochen?
  \item Blut?
  \item Regel?
\end{itemize}
}{}{}{}

\subsection{Der schmerzende Bauch}

\Layout{0}{Der Bauch -- unendliche Weiten}
{%
Gleich vorweg: In vielen
Fällen wird der Grund des Bauchschmerzes nicht mit
hinreichender Sicherheit feststellbar sein und es bietet
sich keine sichere Verdachtsdiagnose an. Statt einer
Diagnose wird dann das \EE{Beschwerdebild beschrieben} (Ort
des Schmerzes, Art des Schmerzes, Dauer, Verlauf,
Austrahlung, Ergebnis der abdominellen Tastuntersuchung:
Abdomen weich oder hart,
Druckschmerz, Loslassschmerz).
}{
\begin{itemize}
  \item In vielen Fällen Krankheitsbild unklar
  \item Im Zweifel: Beschreibung statt Diagnose
\end{itemize}
}{}{}{}

\Layout{2}{Beispiele}
{
\begin{quote}
{\sffamily\bfseries Diagnose:} {\ttfamily Unklarer
Bauchschmerz: Kolikartige Schmerzen im linken Unterbauch seit 1/2h }

{\sffamily\bfseries Anmerkungen:} {\ttfamily keine
Ausstrahlung, Druckschmerz im li. Unterbauch, Abdomen sonst 
weich.}
\vspace{0.5em}

{\sffamily\bfseries Diagnose:} {\ttfamily Unklarer
Bauchschmerz: Brennend-stechender Schmerz im mittleren
Oberbauch seit 2h mit Ausstrahlung in den Rücken}

{\sffamily\bfseries Anmerkungen:} {\ttfamily Abdomen weich,
keine Druck- oder Loslassschmerzen}
\vspace{0.5em}

{\sffamily\bfseries Diagnose:} {\ttfamily Unklarer
Bauchschmerz: „Dumpfer“ Schmerz im linken und rechten
Oberbauch seit gestern. }

{\sffamily\bfseries Anmerkungen:} {\ttfamily Abdomen
weich, keine Druck- oder Loslassschmerzen}
\end{quote}
}{

}{}{}{}


Dennoch muss man bei jedem Patienten versuchen, so gut als
möglich die in Frage kommenden Krankheitsbilder einzugrenzen
und gefährliche Erkrankungen nach Möglichkeit auszuschließen. Im
folgenden sind einige häufige Krankheitsbilder beschrieben,
die \E{erkennbar}  und relativ häufig sind.

\MassnahmenMK{Allgemeine Maßnahmen}{Abdominalerkrankung}
{m:am-abdominalerkrankungen}
{
}{}
% \MassnahmenNG{Allgemeine Maßnahmen}{Abdominalerkrankung}
% {\label{m:am-abdominalerkrankungen}}
% {
% ~
% 
% \begin{itemize}
%   \item \EE{Vitale Bedrohung einschätzen.} \TxBeiVit
%   \TxMassVit
%   \item Lagerung: grundsätzlich \EE{bauchdeckenentspannend}
%   (Knierolle, Beine angezogen), oder nach Wunsch des Patienten
%   \item \EE{Nüchtern lassen} (in Hinblick auf möglicherweise
%   bevorstehende Operation, nichts essen oder trinken lassen)
%   %     \item Monitoring. Der Zustand kann sich abrupt
%   %     verschlechtern.
%   \item \EE{Transportentscheidung}: Situationsgerecht,
%   z.\,B. Abt. f. Chirurgie oder Innere Medizin.
%   
%   Bei Kindern Abt. f. Kinderchirurgie oder Abt. f.
%   Kinderheilkunde (je nach Verdachtsdiagnose und regionaler
%   Regelung).
%   
%   Bei weiblichen Patienten Abt. für Gynäkologie (evtl. im
%   \enquote{Hintergrund}) erwägen.
% \end{itemize}
% }


% \clearpage % TEMP temp clearpage
\section{Häufige und gut erkennbare Erkrankungen}



\subsection{Das Akute Abdomen}
\label{topic:akutes-abdomen}
\index{Abdomen!akutes|see{Akutes Abdomen}}


\Definition{\T{Akutes Abdomen}: Akute
Erkrankung des Bauchraumes mit plötzlichen, starken
\EE{Bauchschmerzen} und
\EE{bretthartem Bauch}, welche im weiteren Verlauf
lebensbedrohlich ist. Es sind viele Ursachen möglich!}

\Commentary[Akutes Abdomen, Definition]{Die Definition des
Akuten Abdomens ist oft nicht ganz eindeutig: Einerseits kann man 
darunter jede plötzlich auftretende
Baucherkrankung
verstehen \citep{Longmore2007,Renz-Polster2006}. 

Allerdings ver\-stehen viele Leute darunter eher eine
bedrohliche Bauch\-erkrankung, bei der man akut (schnell)
handeln muss (Ab\-wehrspannung, harter Bauch). In der
Notfallmedizin halten wir uns an letztere Definition.}


\Layout{2}{\TxBeschreibung}
{Das \E{Akute Abdomen} ist eine Kombination von bestimmten
Symptomen, die
durch viele verschiedene Erkrankungen im Bauchraum ausgelöst
werden kann. Diese zu Grunde liegenden Erkrankungen können
plötzlich entstehen, oder sich über längere Zeit entwickeln
und erst im fortgeschrittenen Stadium die Symptome des
akuten Abdomens auslösen. Ein akutes Abdomen erfordert eine
rasche (spitalsmäßige) Behandlung.
}
{

}{}{}{}

\Layout{0}{\TxSymptome}
{%
Der Patient liegt meistens im Bett. Oft ist eine
Schonhaltung (angezogene Beine, gekrümmte Körperhaltung) zu
beobachten. Er klagt über \EE{Schmerzen} im Bauch, diese können
als diffus (im gesamten Bauch) oder auf ein Bauchsegment
beschränkt angegeben werden.

Die \EE{harte Bauchdecke} (\enquote{\E{bretthart}}) ist
neben den Schmerzen das Leitsymptom des Akuten Abdomens.
Es besteht oft Übelkeit oder auch Erbrechen. Der Kreislauf
kann kompensiert sein, oder Zeichen von Schwäche zeigen, bis
hin zu einer Schocksymptomatik (Blutverlust, Sepsis).
Weitere Symptome sind abhängig von der zu Grunde liegenden
Erkrankung.

Diese Auflistung der Symptome ist in der Praxis nicht
unbedingt vollständig bei einem Patienten nachzuweisen. Die
Intensität und die Anzahl der angeführten Symptome ist
einerseits vom Allgemeinzustand, vom Lebensalter, von
anderen vorbestehenden Erkrankungen des Patienten und der
Dauer des aktuellen Zustandes abhängig.

Der Patient ist nicht unbedingt unmittelbar vital
bedroht. Wesentlich nach Eintreffen beim Patienten ist die
Beantwortung der Fragen \Speech{\enquote{Wie schlecht, bzw.
wie gut geht es dem Patienten? Ist der Patient \uline{akut}
vital bedroht?}}.

} 
{
\begin{itemize}
  \item \EE{Schmerz}: diffus (ganzer Bauch) oder
  lokalisierbar
  \item \EE{Harte  Bauchdecke} (bretthart)
  \item Schonhaltung  (angezogene Beine, Patient
  krümmt sich)
  \item Übelkeit, Erbrechen
  \item Evtl. Stuhl-, Windverhalten
  \item Evtl. Fieber
  \item Evtl. Schock
\end{itemize}
}{}{}{}




\Layout{0}{Ursachen}
{%
Grundsätzlich kommt fast jede Erkrankung oder Verletzung in
Betracht die Strukturen und Organe im Bauchraum betrifft,
entsprechend vielfältig sind auch die Ursachen: 

\begin{multicols}{2}
\FormatTable
\begin{itemize}
  \item Bauchfellentzündung (Peritonitis)
  \item Darmverschluss (Ileus)
  \item Blinddarmentzündung (Appendizitis)
  \item Entzündung der Bauchspeicheldrüse  (Pankreatitis)
  \item Entzündung des Magens / Magengeschwür
  \item Durchbruch / Zerreißung eines Hohlorganes
  \item Blutungen in den Bauchraum
  \item Bauchaortenaneurysma
  \item Mesenterialinfarkt
  \item andere abdominelle Erkrankungen
  \item Entzündung der Leber (Hepatitis )
  \item Gynäkologische Erkrankungen
  \item \ldots
\end{itemize}
\end{multicols}
}{
\begin{itemize}
  \item Bauchfellentzündung 
  \item Blinddarmentzündung 
  \item Entzündung der Bauchspeicheldrüse  (Pankreatitis)
  \item Durchbruch eines Hohlorganes
  \item Blutungen in den Bauchraum
  \item Bauchaortenaneurysma
  \item Gynäkologische Erkrankungen
  \item u.\,v.\,a.\,m. \ldots
\end{itemize}
}{}{}{}



\MassnahmenMK{Spezielle Maßnahmen}{Akutes Abdomen}
{m:akutes-abdomen}{
\EE{Taktik}: \Speech{Unmittelbare vitale Bedrohung
einschätzen, sie ist
abhängig von den Vitalwerten.
Vermutliche Ursache und Differentialdiagnosen
abklären, Transport an eine geeignete Abteilung}
}{}

% \MassnahmenNG{Spezielle Maßnahmen}{Akutes Abdomen}
% {\label{m:akutes-abdomen}}{
% \EE{Taktik}: \Speech{Unmittelbare vitale Bedrohung
% einschätzen, sie ist
% abhängig von den Vitalwerten.
% Vermutliche Ursache und Differentialdiagnosen
% abklären, Transport an eine geeignete Abteilung}
% 
% 
% \begin{itemize}
%   \item Allgemeine Maßnahmen bei Abdominalerkrankungen
%   (\prettyref{m:am-abdominalerkrankungen})
%   \item \EE{Vitale Bedrohung einschätzen.} Ggfs.
%   \TxMassVit
%   \item Lagerung: grundsätzlich \EE{bauchdeckenentspannend}
%   (Knierolle, Beine angezogen), oder nach Wunsch des Patienten
%   \item \EE{Nüchtern lassen} (in Hinblick auf möglicherweise
%   bevorstehende Operation, nichts essen oder trinken lassen)
%   \item Monitoring: Der Zustand kann sich abrupt
%   verschlechtern.
%   \item \EE{Transportentscheidung}: Abt. f. \EE{Chirurgie},
%   bei Kindern ${\rightarrow}$ Abt. f. \EE{Kinderheilkunde},
%   bei Frauen evtl. \EE{Gynäkologie} oder Chirurgie mit Abt.
%   f. Gynäkologie im Hintergrund
% \end{itemize}
% }


\subsection{Darmverschluss}
\index{Darmverschluss}
\label{topic:ileus}
\HeadingSub{\Lang{Term.} \T{Ileus}
(vollständig), \Z{\T{Subileus}  (unvollständig)}}

\Definition{Unterbrechung der normalen  Darmpassage}

\Layout{2}{\TxBeschreibung}
{%
Der Verschluss kann mechanisch bedingt sein (durch
Hindernisse wie Verwachsungen oder Tumore), oder durch eine
Lähmung der Darmmuskulatur (Medikamente, Gifte, Entzündung
\ldots) verursacht sein. Bestimmte Schmerzmittel
\Z{(Opiate)}, welche oft bei Krebspatienten eingesetzt
werden, sind bekannt dafür Darmlähmungen zu erzeugen. 
}{%
\begin{itemize}
  \item Mechanische Blockade oder Lähmung
\end{itemize}
}{}{}{}


\Layout{0}{\TxSymptome}
{%
Typisch für den Darmverschluss ist die mangelhafte
Stuhlproduktion. Es kommt zu diffusen Bauchschmerzen und einem
Blähbauch, sowie zu Übelkeit und Erbrechen. In seltenen
Fällen kann es auch zum Erbrechen von zurück
gestautem Kot kommen \Z{(\I[Miserere|Z]{Miserere})}.

\bigskip

\SlideCompensationNoParskip{1}
}{
\begin{itemize}
  \item Diffuse Bauchschmerzen
  \item Stuhlverhalten
  \item Übelkeit, Erbrechen (selten: Koterbrechen)
\end{itemize}
}{}{}{}


\MassnahmenMK{Spezielle
Maßnahmen}{Darmverschluß}{m:darmverschluss}{
\EE{Taktik}: \Speech{Unmittelbare Bedrohung je nach
Vitalwerten. Nüchtern lassen, Differentialdiagnosen
abklären, Transport an geeignete Einrichtung}
}{}
% \MassnahmenNG{Spezielle
% Maßnahmen}{Darmverschluß}{\label{m:darmverschluss}}{
% \EE{Taktik}: \Speech{Unmittelbare Bedrohung je nach Vitalwerten. Nüchtern lassen, Differentialdiagnosen
% abklären, Transport an geeignete Einrichtung}
% 
% 
% \begin{itemize}
%   \item Allgemeine Maßnahmen bei Abdominalerkrankungen
%   (\prettyref{m:am-abdominalerkrankungen})
%   \item \EE{Vitale Bedrohung einschätzen}, ggfs.
%   \TxMassVitBes
%   \begin{itemize}
%     \item Lagerung: bauchdeckenentspannend, bzw. je nach
%     Patientenwunsch
%   \end{itemize}
%   \item Patient nüchtern lassen.
%   \item Transportentscheidung: Abt. f. Chirurgie
% \end{itemize}
% }


\subsection{Blinddarmentzündung}
\label{topic:appendizitis}
\HeadingSub{\Lang{Term.} \T{Appendizitis}}


\Definition{Entzündung des \EE{Wurmfortsatzes}
(\EE{Appendix} \Z{vermiformis}) des Blinddarms}

\Layout{0}{\TxSymptome{} und Verlauf}
{%
\ACh{GABG}{2009-05-07}{NEU}Im Vordergrund der akuten Appendizitis
steht anfänglich meist ein als eher diffus empfundener,
heftiger Schmerzzustand im Bereich der
Nabelgegend\ACh{GABS}{2009-05-07}{Magengegend --> Nabelgegend}, der
innerhalb weniger Stunden in den rechten Unterbauch
wandert. Übelkeit, Brechreiz und Verstopfung
\ACh{GABS}{2009-05-07}{Obstipation --> Verstopfung} können
das Krankheitsbild begleiten. Meist ist die Körpertemperatur
bis 38\DegC{} leicht erhöht
\ACh{GABS}{2009-05-07}{gelöscht da unerheblich für RS: , wobei eine rektale und
axilläre Messung eine Differenz um ca. 1\DegC aufweist}. 
\ACh{GABS}{2009-05-07}{alt: Unbehandelt kann der weitere Verlauf durch eine
Abszeßbildung und Perforation in
die Bauchhöhle gekennzeichnet sein. neu: Unbehandelt kann es zu einer Eiterung
und einem Durchbruch in die Bauchhöhle kommen.} Unbehandelt kann es zu einer
Eiterung und einem Durchbruch in die Bauchhöhle, einer
anschließenden Bauchfellentzündung und einem akuten Abdomen
kommen (\FullPrettyRef{topic:akutes-abdomen}).

Der Verlauf einer akuten Appendizitis ist jeweils durch
individuelle Faktoren wie Lebensalter und Allgemeinzustand
des Patienten geprägt.
}{
\begin{itemize}
  \item Schmerzen im re. Unterbauch
  \item \E{Druck}schmerzen im re. Unterbauch
  \item \E{Loslass}schmerz li. Bauch
  \item Abwehrspannung
  \item Übelkeit, Erbrechen
  \item Evtl. Fieber
\end{itemize}
}{}{}{}



\MassnahmenMK{Spezielle
Maßnahmen}{Appendizitis}{m:appendizitis}{
\EE{Taktik}:
\Speech{Differntialdiagnosen abklären, Transport an geeignete Abteilung}
}{}
% \MassnahmenNG{Spezielle
% Maßnahmen}{Appendizitis}{\label{m:appendizitis}}{
% \EE{Taktik}:
% \Speech{Differntialdiagnosen abklären, Transport an geeignete Abteilung}
% 
% 
% \begin{itemize}
%   \item Allgemeine Maßnahmen bei Abdominalerkrankungen
%   (\prettyref{m:am-abdominalerkrankungen})
%   \item Patient nüchtern lassen.
%   \item Transportentscheidung: Abt. f. Chirurgie, Sonderfälle
%   (Kinder, Frauen!) beachten!
% \end{itemize}
% }



\subsection{Gallenkolik}
\label{topic:gallenkolik}
\index{Kolik!Gallen-}
\index{Gallenkolik}

\Definition{Kolik-artige Schmerzen infolge einer Blockade
der Gallenwege durch Gallensteine}

\Z{\emph{Vgl. Kolik: \prettyref{topic:kolik}}}

\Layout{0}{\TxBeschreibung} 
{Gallensteine verursachen eine Verstopfung des
Gallengangs und folglich einen Rückstau von
Gallenflüssigkeit in die Leber. Dabei wird der Gallengang
gedehnt und es kommt zu starken Schmerzen.} 
{
\begin{itemize}
  \item Verstopfung der Gallenwege (Stein)
  \item → Rückstau und Dehnung
\end{itemize}
}{}{}{}

\SlideCompensationNoParskip{2}

\Layout{0}{\TxSymptome}
{%
Durch den Rückstau und der Dehnung des Gallengangs kommt es
zu starken, auf- und abschwellenden, krampfartigen
(\EE{kolikartigen}) \EE{Schmerzen} im \EE{rechten
Oberbauch}, klassischerweise mit Ausstrahlung in die rechte
Schulterregion. Begleitendes Fieber ist möglich.

Der rechte Oberbauch ist äußerst druckschmerzhaft, der
Patient ist sehr \EE{unruhig} und agitiert.

\bigskip

\bigskip

\bigskip
% \SlideCompensationNoParskip{1}
}{
\begin{itemize}
  \item Auf- und abschwellende, krampfartige
  (kolikartige) Schmerzen im re. Oberbauch
  \item Nach fettreichen Mahlzeiten
  \item Evtl. Gelbfärbung des Augenweiß (Gelbsucht)
  \item Evtl. Übelkeit, Erbrechen
\end{itemize}
}{}{}{}

\MassnahmenMK{Spezielle
Maßnahmen}{Gallenkolik}{m:gallenkolik}{
\EE{Taktik}: \Speech{Rasche Schmerztherapie (veranlassen)}
}{}
% \MassnahmenNG{Spezielle
% Maßnahmen}{Gallenkolik}{\label{m:gallenkolik}}{
% \EE{Taktik}: \Speech{Rasche Schmerztherapie (veranlassen)}
% 
% 
% \begin{itemize}
%   \item Allgemeine Maßnahmen bei Abdominalerkrankungen
%   (\prettyref{m:am-abdominalerkrankungen})
%   \item Rasche Schmerztherapie veranlassen: Je nach
%   Situation, Notarzt oder rascher Transport in eine
%   geeignete Einrichtung
%   \item Patient nüchtern lassen
%   \item Transportentscheidung: Abt. f. Chirurgie
% \end{itemize}
% }


% \clearpage % TEMP temp clearpage
\subsection{Gastrointestinale
Blutungen und Ösophagusvarizenblutung}
\label{topic:oesophagusvarizenblutung}
\label{topic:gastrointestinale-blutung}
\TODO{EMHO}{2010-08-04}{andere Blutungsquellen sind
anzuführen}{}

\Definition{Blutungen im Magen-Darm-Trakt bzw. blutende \enquote{Krampfadern} in der
Speiseröhre.}\nopagebreak[4]

\Layout{0}{\TxBeschreibung}
{Im Verlauf des Magen-Darm-Traktes kann es an verschiedenen
Stellen zu Blutungen kommen, die auch sehr schwer ausfallen
können. 

Blutungen im \EE{oberen Teil} des Verdauungstraktes können
durch \T{Ösophagusvarizen} verursacht werden: Bei Patienten
mit einer fortgeschrittenen Lebererkrankung (Leberzirrhose, \ref{topic:leberzirrhose}) kann es zur Bildung von Krampfader-artigen Gefäßerweiterungen in der Speiseröhre
kommen. Diese Gefäße können massive Blutungen verursachen,
die tödlich enden können. Auch Blutungen im Rahmen eines
Magengeschwürs sind möglich, aber selten dermaßen stark
ausgeprägt.

\Z{Blutungen im \EE{mittleren Teil} des Magen-Darm-Traktes
können durch diverse Erkrankungen und Tumore verursacht
werden. Blutungen im \EE{Endteil} kommen oft durch Tumore
oder infolge eines Hämorrhoidalleidens zu Stande.}


}
{
\begin{itemize}
    \item Oberer Teil: Ösophagusvarizzen
    
    (Oft bei
    Lebererkrankungen, hoher Blutverlust möglich)
    \item \Z{Diverse Erkrankungen, Tumore}
\end{itemize}
}{}{}{}

\Layout{0}{\TxSymptome}
{
Bei Blutungen im \EE{oberen Magen-Darm-Trakt}
(Speiseröhre, Magen) kommt es je nach Stärke zu
schwallartigem Erbrechen  von rotem, frischen Blut
(\Z{\I{Hämatemesis}}) oder häufiger zum Erbrechen von
\II[Erbrechen!Kaffeesatz-artiges]{Kaffeesatz-artigen}
Mageninhalt: Durch die Magensäure kommt es zur Veränderung
des Blutes, es wird bräunlich, ähnlich dem Kaffeesatz im
Kaffeefilter. Bei Blutungen im \EE{mittleren Teil} des
Verdauungstraktes färbt sich der Stuhl tief schwarz, man
spricht vom sog. \T{Teerstuhl} (\T{Melena}. Bei Butungen im
\EE{Endteil} wird rötliches Blut ausgeschieden. 

}{
\begin{itemize}
  \item Oberer Teil: Erbrechen von
  \begin{itemize}
    \item Blut im Schwall
    \item Kaffeesatz-artiger Mageninhalt
  \end{itemize} 
  \item Mittlerer Teil: Teerstuhl
  \item Endteil: Blutiger Stuhl
\end{itemize}
}{}{}{}

\textleaf


\Fazit{Der Patient kann aufgrund des Blutverlustes akut
vital bedroht sein, es besteht die Gefahr eines
hypovolämischen Schocks.}

\MassnahmenMK{Spezielle
Maßnahmen}{Blutungen
des Verdauungstrakts}{m:blutung-verdauungstrakt}{
\EE{Taktik}: \Speech{\Ca{Vitale Bedrohung bei hohem Blutverlust!}
Schockbekämpfung, zügiger Transport}
}
% \MassnahmenNG{Spezielle
% Maßnahmen}{Blutungen
% des Verdauungstrakts}{\label{m:blutung-verdauungstrakt}}{
% \EE{Taktik}: \Speech{\Ca{Vitale Bedrohung bei hohem Blutverlust!}
% Schockbekämpfung, zügiger Transport}
% 
% 
% \begin{itemize}
%   \item \TxMassVit
%     \item Lagerung:
%     Situationsgerecht, je nach Blutungsquelle und
%     Patientenzustand
%   \item Allgemeine Maßnahmen der Schockbekämpfung
%   (\prettyref{m:am-schock})
%   \item Ggfs. spezielle Maßnahmen bei hypovolämischem Schock
%   (\prettyref{topic:schock-hypovolaemischer})
%   \item Transportentscheidung: Abteilung für Chirurgie,
%   Voranmeldung!
% \end{itemize}
% }






\subsection{Magen-Darm-Grippe, Gastroenteritis,
Lebensmittelvergiftung}
\label{topic:gastroenteritis}
\label{topic:lebensmittelvergiftung}


\HeadingSub{\Lang{Syn.} Brechdurchfall,
\Z{gastrointestinaler Infekt, Gastroduodenitis}}

Die Magen-Darm-Grippe ist eine chronisch oder akut
verlaufende Entzündung der Magen- und der Dünndarmschleimhaut. 

\Layout{0}{\TxBeschreibung{} und Ursachen}
{%
Die Gastroenteritis wird durch Erreger wie Bakterien oder
Viren hervorgerufen. Es gibt völlig harmlose aber auch
lebensbedrohliche Verlaufsformen, dies ist abhängig vom
Erreger und letztendlich auch vom Alter des Patienten. Durch
den Durchfall und das Erbrechen kann es zu einem maßgeblichen
Verlust von Wasser und Elektrolyten kommen. Wenn der Patient
nicht ausreichend trinken kann besteht die Gefahr der
Austrocknung (\I[Exsikkose!bei gostrointestinalem
Infekt]{Exsikkose}). 
}{
\begin{itemize}
  \item Vielfältig
  \item Häufig Infektionen
\end{itemize}
}{}{}{}


\Layout{2}{\TxSymptome}
{
Klassische Zeichen einer Magen-Darm-Grippe sind Erbrechen
(\Z{\I{Emesis}}) und Durchfall (\Z{\I{Diarrhoe}}). Dazu hat
der Patient meist starke, kaum lokalisierbare
Bauchschmerzen. Infolge des 
Erbrechens und des Durchfalls kann es zu einem erheblichen 
Wasser- und Elektrolytverlust und Symptomen der Exsikkose
(\ref{topic:exsikkose}) kommen. 

}{
\begin{itemize}
    \item Erbrechen 
    \item Durchfall 
    \item Bauchschmerzen
    \item Exsikkosezeichen
\end{itemize}
}{}{}{}

\Layout{0}{Besonders gefährdet}
{%
Besonders betroffen davon sind Kleinkinder, da sie
normalerweise einen höheren Wasseranteil am Körpergewicht
haben.  Die Symptome sind klassisch: Oberbauchkrämpfe,
die ziemlich heftig sein können, Übelkeit, Erbrechen, Durchfälle.
Exsikkose-Gefahr! 
}{
\begin{itemize}
    \item (Klein-)Kinder
    \item Alte Menschen
\end{itemize}
}{}{}{}


\MassnahmenMK{Spezielle
Maßnahmen}{Gastroenteritis}{m:gastroenteritis}{
}{}
% \MassnahmenNG{Spezielle
% Maßnahmen}{Gastroenteritis}{\label{m:gastroenteritis}}{
% \begin{itemize}
%   \item Allgemeine Maßnahmen bei Abdominalerkrankungen
%   (\prettyref{m:am-abdominalerkrankungen})
%   \item Patient nüchtern lassen
%   \item Transportentscheidung: Abt. f. Innere Medizin; bei
%   Kindern Abt. f. Kinderheilkunde
%   \item Strikte Desinfektionsmaßnahmen!
% \end{itemize}
% }







\subsection{Flüssigkeitsmangel: Exsikkose}
\ACh{WITM}{????-??-??}{NEU: Schlagworte}
\ACh{GABG}{2009-05-11}{NEU: Fließtexte}
\ACh{GABS}{2009-05-14}{Fließtexte stark modifiziert und angepasst. }
\label{topic:exsikkose}

\HeadingSub{\Lang{Syn.} \T{Dehydratation}}

\Definition{\enquote{Austrocknung}: Es besteht  ein
hochgradiges Defizit von Flüssigkeit im Körper.}

Grundsätzlich wird (wurde) entweder \E{zu wenig Flüssigkeit
zugeführt} oder \E{zu viel Flüssigkeit an die Außenwelt
abgegeben} (eine Kombination aus beidem ist natürlich auch
möglich). Dafür gibt es jeweils sehr verschiedene Ursachen:

\Layout{0}{Ursachen}
{%
Ältere Menschen haben oft ein mangelndes Durstgefühl und daher
einen verminderten Trinkantrieb. Auch durch die Einnahme von
Drogen (Ecstacy) kann auf den Durst \enquote{vergessen} werden.  

Andererseits kann es durch \E{Erbrechen}, \E{Durchfall}
sowie \E{starkes Schwitzen} zu einem stark erhöhten
Flüssigkeitsverlust kommen. Bei \E{Fieber} wird auch
vermehrt Feuchtigkeit abgeatmet, und natürlich können
auch \E{Medikamente} \enquote{entwässern}.
}{
\begin{itemize} %GABS: umgebaut
  \item Flüssigkeitszufuhr ${\downarrow}$
  \begin{compactitem}
    \item Verminderter Trinkantrieb
  \end{compactitem}
  \item Flüssigkeitsverlust ${\uparrow}$
  \begin{compactitem}
    \item Erbrechen, Durchfall
    \item Starkes Schwitzen
    \item Fieber (Feuchtigkeit wird vermehrt abgeatmet)
    \item Medikamente
  \end{compactitem}
\end{itemize}
}{}{}{}

Die genannten Ursachen sind natürlich nur ein kleiner Teil
von sehr vielen Möglichkeiten.

\Layout{0}{\TxSymptome}
{%
Der Patient wirkt allgemein \E{geschwächt bis apathisch}, oft
\EE{verwirrt}, klagt über mehr oder minder heftiges
Durstgefühl, ev. über Wadenkrämpfe und Erbrechen. Die Haut
ist auffallend trocken, \EE{Falten sind abhebbar und bleiben
bestehen}. \EE{Mundhöhle und Zunge sind trocken}, evtl.
borkig oder mit eingetrocknetem Schleim belegt. Die Augen
sind charakteristischerweise tiefliegend, die Nase wirkt
auffallend spitz. Der RR ist mitunter sehr nieder, die HF
eher hoch. \E{Der Patient ist bis zur Klärung der Situation
nüchtern zu belassen!}
}{
\begin{itemize}
    \item RR ${\downarrow}$
    \item stehende Hautfalten
    \item trockene, belegte Zunge
    \item Elektrolytentgleisung
    \item \EE{Verwirrtheit, Schwindel, Apathie}
\end{itemize}
}{}{}{}



% \clearpage % TEMP temp clearpage
\section{Weniger häufige oder nicht-so-gut-erkennbare Erkrankungen}

% \minisec{Querverweise}
% 
% \begin{itemize}
%     \item Hepatitis: \prettyref{topic:hepatitis}
% \end{itemize}




\subsection{Magen- und Zwölffingerdarmgeschwür}
\label{topic:gastritis}

\Z{\HeadingSub{\Lang{Synn.} Gastritis, Duodenitis}}

\Definition{Magengeschwür: Entzündung der
Magenschleimhaut mit Substanzdefekt.}

\Layout{0}{Ursachen}
{
Die häufigste Ursache für ein Magengeschwür ist eine
chronische Infektion mit einem Bakterium \Z{(Helicobacter
pylori)} \citep{Malfertheiner2004}. Medikamente, wie
z.\,B. nicht-steroidale Antirheumatika (NSAR, z.\,B.
Voltaren{\texttrademark} (Diclofenac),
Aspirin{\texttrademark} (Acetylsalicylsäure, ASS),etc.)
schädigen die Magenschleimhaut und können ebenso bei
häufiger Einnahme zu einem Geschwür führen. Auch andere
Stoffe können auf chemisch-toxischen Weg zu einer
entsprechenden Schädigung führen. \Z{Ein
Gallensäurenrückfluß (Reflux) bzw. Autoimmunerkrankungen
können ebenso zu einem Geschwür führen.} 
}{
\begin{itemize}
    \item Bakteriell \Z{(Helicobacter pylori)}
    \item Medikamente
    \item Chemisch-toxisch
    \item \Z{Gallensäurenrückfluß}
    \item \Z{Autoimmunerkrankungen}
\end{itemize}
}{}{}{}

\Layout{0}{\TxSymptome}
{Der Patient klagt über einen
\EE{stechend-schneidenden Schmerz}, welcher sehr stark sein
kann. Daneben kommt es zu Appetitlosigkeit, Völlegefühl,
Übelkeit und Erbrechen. Wenn gleichzeitig eine \EE{Blutung}
besteht, dann sind \II[Teerstuhl!Gastritis]{Teerstühle}
(schwarz, übelriechend) oder \EE{Kaffeesatz-artiges
Erbrechen} zu beobachten. 
}{
\begin{itemize}
    \item Schmerzen (Oberbauch), Übelkeit, Erbrechen
    \item Kaffeesatz-artiges Erbrechen, Teerstuhl
\end{itemize}
}{}{}{}


\MassnahmenMK{Spezielle Maßnahmen}{Gastritis,
Duodenitis}{m:gastro-duodenitis}{
}{}
% \MassnahmenNG{Spezielle Maßnahmen}{Gastritis,
% Duodenitis}{\label{m:gastro-duodenitis}}{
% \begin{itemize}
%   \item Allgemeine Maßnahmen bei Abdominalerkrankungen
%   (\prettyref{m:am-abdominalerkrankungen})
%   \item \EE{Transportentscheidung}: 
%   Abt. f. Innere Medizin 
%   
%   Bei Blutungszeichen (Teerstuhl, kaffeesatzartiges
%   Erbrechen): Abt. f. Chirurgie
% \end{itemize}
% }

\subsection{Pankreatitis}

\HeadingSub{\Lang{Syn.} \T{Bauchspeicheldrüsenentzündung}}
\label{topic:pankreatitis}

Wie nahezu jedes Organ des Körpers kann sich auch die
Bauchspeicheldrüse entzünden. Die
Bauchspeichedrüsenentzündung (Pankreatitis) kann dabei
sowohl \E{akut}, als auch \E{chronisch} verlaufen. 

% \Z{\Lang{ICD-10} Akute Pankreatitis: K85, chronisch
% Pankreatitis: K86.1}

\Layout{0}{Ursachen}
{\index{Gallensteine!Pankreatitis}
\index{Alkoholmißbrauch!Pankreatitis}
Die häufigsten Ursachen für eine Pankreatitis sind
\EE{Gallensteinleiden} und \EE{Alkoholmißbrauch}
\citep{SiewertChirurgie8-07.14-Pankreas,ScholzNotfallmedizin2}.
Daneben gibt es noch viele, allerdings nicht so häufige
Ursachen, oft bleibt die Ursache auch unbekannt
\citep{HeroldInnereMedizin2005}.

\Speech{Warum können Gallensteine eine Pankreatitis
verursachen?} Bei vielen Menschen mündet der
Pankreasausführungsgang, welcher Verdauungsenzyme in de
Zwölffingerdarm abgibt, nicht direkt in den Dünndarm,
sondern in den unteren Abschnitt des Gallengangs.
Verschließt nun ein Stein den Gallengang unterhalb der
Einmündung, so staut sich Galle und Pankreassaft in die
Bauchspeicheldrüse zurück. Es kann sogar dazu kommen, dass
sich das Pankreas durch die Verdauungsenzyme selber
an-/verdaut.

}
{\begin{itemize}
  \item Gallensteinleiden
  \item Alkoholmißbrauch
  \item \ldots
\end{itemize}}
{}
{}
{}

\Layout{0}{\TxSymptome}
{Die Symptome sind oft unspezifisch, eine
Erstdiagnosestellung ist präklinisch normalerweise nicht
möglich. Charakteristisch ist bestenfalls ein
\II{gürtelförmiger
Oberbauchschmerz}\index{Oberbauchschmerz,gürtelförmiger}.
Die häufigsten Begleitsymptome sind Übelkeit und Erbrechen,
sowie Fieber. \citep{HeroldInnereMedizin2005}


\bigskip

\bigskip
}
{\begin{itemize}
  \item Oberbauchschmerzen, oft gürtelförmig
  \item Übelkeit, Erbrechen
  \item Fieber
\end{itemize}}
{}
{}
{}

\Layout{0}{\TxKomplikationen}
{Je nach Schwere und Dauer der Entzündung kann es zum
Selbstverdauen der Bauchspeicheldrüse bzw. zu einer
\E{Ruptur} kommen. In Folge tritt ein
\I[Akutes Abdomen!bei Pankreatitis]{akutes Abdomen} auf
(\FullPrettyRef{topic:akutes-abdomen}) und der Patient kann
kann einen \I[Schock!septischer!bei Pankreatitis]{septischen
Schock} erleiden (\FullPrettyRef{topic:schock-septischer}).

\bigskip

}
{\begin{itemize}
  \item Ruptur
  \item Akutes Abdomen
  \item Septischer Schock
\end{itemize}}
{}
{}
{}

\bigskip

\MassnahmenMK{Spezielle Maßnahmen}{Akute
Pankreatitis}{m:pankreatitis-akut}{
}{}
% \MassnahmenNG{Spezielle Maßnahmen}{Akute
% Pankreatitis}{\label{m:pankreatitis-akut}}{
% \begin{itemize}
%   \item Allgemeine Maßnahmen bei Abdominalerkrankungen
%   (\prettyref{m:am-abdominalerkrankungen})
%   \item \EE{Transportentscheidung}: Abt. f. Chirurgie
% \end{itemize}
% }


\subsection{Mesenterialinfarkt}
\label{topic:mesenterialinfarkt}

Beim Mesenterialinfarkt kommt es zum Verschluss einer den
Darm versorgenden Arterie
\Z{(\I[Mesenterialarterie|Z]{Mesenterialarterie}\ZFootnote{Das
\I[Mesenterium!Mesenterialinfarkt]{Mesenterium}
(\I{Darmgekröse}) ist die Aufhängung von Teilen des
Dünndarms (\FullPrettyRef{topic:mesenterium}). Darin
verlaufen venöse und arterielle Gefäße zum und vom Darm.})}
mit einem Absterben des entsprechenden Darmabschnittes. Er
ist somit ein \I[Ischämie!Mesenterialinfarkt]{ischämischer}
Notfall.

\Layout{0}{\TxSymptome}
{Die Symptome sind eher unspezifisch. Der Patient klagt oft
(aber nicht immer) über \EE{abdominelle Schmerzen}, welche
oft dumpf und diffus sind. Es kann für einige Stunden ein
\E{stilles Intervall} folgen, in welchem die Scherzen eher
gering sind. 

Nach ungefähr \VonBis{12}{24}\Hour* wird ein
\EE{Darmverschluss} bemerkbar und die Darmwand verliert
zunehmend an Dichtigkeit, es kann zu einem \EE{Durchbruch}
(Ruptur) kommen. Es folgt eine \EE{Bauchfellentzündung} und
ein \II[Akutes Abdomen!beim Mesenterialinfarkt]{akutes
Abdomen} (\FullPrettyRef{topic:akutes-abdomen}). Es kann zu
\EE{blutigen, übel riechenden Stuhlabgängen} kommen.
\cite{SiewertChirurgie8-06-Gefaesschirurgie}
% 
% \bigskip

}
{\begin{itemize}
  \item Adbominalschemrzen
  \item Übelkeit, Erbrechen
  \item $\sim$ 6\Hour* stilles Intervall
  \item Nach $\sim$ \VonBis{12}{24}\Hour* Durchbruch und
  Bauchfellentzündung, Ileus
  \item Übelriechende, blutige Stuhlabgänge
\end{itemize}}
{}
{}
{}

\MassnahmenMK{Spezielle
Maßnahmen}{Mesenterialinfarkt}{m:mesenterialinfarkt}{
}{}
% \MassnahmenNG{Spezielle
% Maßnahmen}{Mesenterialinfarkt}{\label{m:mesenterialinfarkt}}{
% \begin{itemize}
%   \item Allgemeine Maßnahmen bei Abdominalerkrankungen
%   (\prettyref{m:am-abdominalerkrankungen})
%   \item Nüchtern lassen!
%   \item \EE{Transportentscheidung}: Abt. f.
%   Chirurgie
% \end{itemize}
% }

% \clearpage % TEMP temp clearpage
\subsection{Bauchfellentzündung}
\label{topic:bauchfellentzuendung}
\label{topic:peritonitis}

\HeadingSub{\Lang{Term.} \T{Peritonitis}}

\Definition{Entzündung und Keimbesiedelung des 
Bauchfells}

\Layout{2}{\TxBeschreibung}{%
Durch die Infektion und Entzündung des Bauchfelles kommt es
zu Symptomen eines Akuten Abdomens. Oft ist die Perforation
eines Hohlorganes  (Magendurchbruch, Darmperforation,
Blinddarmdurchbruch, \ldots) Ursache der Entzündung. 
}{

}{}{}{}

\Layout{0}{Ursache}{
\begin{itemize}
    \item Perforation eines Hohlorganes (Darm, Magen,
    Gallenblase, \ldots)
    \item Fremdkörper
\end{itemize}
}{
\begin{itemize}
    \item Perforation eines Hohlorganes
    \item Infektion über Blut- / Lymphwege oder eingewandert
    über Fremdkörper
\end{itemize}
}{}{}{}

\Layout{0}{\TxSymptome}{
\begin{itemize}
    \item Symptome eines Akuten Abdomens
    \item Schmerzen (diffus)
    \item Evtl. reflektorischer Darmverschluss
    \item Evtl. Sepsis- und/oder Schockzeichen
\end{itemize}

\bigskip

}{
\begin{itemize}
    \item Symptome eines Akuten Abdomens
    \item Schmerzen (diffus)
    \item Evtl. reflektorischer Darmverschluss
    \item Evtl. Sepsis- und/oder Schockzeichen
\end{itemize}
}{}{}{}


\MassnahmenMK{Spezielle
Maßnahmen}{Bauchfellentzündung}{m:bauchfellentzuendung}{
}{}
% \MassnahmenNG{Spezielle
% Maßnahmen}{Bauchfellentzündung}{\label{m:bauchfellentzuendung}}{
% \begin{itemize}
%   \item Maßnahmen wie bei akutem Abdomen
%   (\FullPrettyRef{m:akutes-abdomen}), und sonst abhängig von
%   der Ursache
% \end{itemize}
% }


% \Z{
% \subsection{Reserviert \A{Sonstige -- nicht akute --
% Erkrankungen}}
% \subsubsection{Reserviert \A{Chronisch entzündliche
% Darmerkrankungen}} }
% 

